\documentclass[12pt,a4paper]{article}
\usepackage[utf8]{inputenc}
\usepackage{amsmath,amssymb,amsthm}
\usepackage{geometry}
\geometry{margin=1in}

\newtheorem{theorem}{Theorem}
\newtheorem{lemma}[theorem]{Lemma}
\newtheorem{corollary}[theorem]{Corollary}
\newtheorem{proposition}[theorem]{Proposition}
\theoremstyle{definition}
\newtheorem{definition}[theorem]{Definition}

\title{Problem 158: Exploring Strings with Exactly One Ascent}
\author{Project Euler Solution}
\date{}

\begin{document}
\maketitle

\section{Problem Statement}

For $1 \le n \le 26$, let $p(n)$ be the number of strings of length $n$ formed from $n$ distinct letters chosen from the 26-letter alphabet, such that exactly one character is lexicographically after its left neighbor. Find $\max_{1 \le n \le 26} p(n)$.

\section{Mathematical Derivation}

\subsection{Eulerian Numbers}

An \emph{ascent} in a permutation $\sigma$ of $\{1, \ldots, n\}$ is a position $i$ where $\sigma(i) < \sigma(i+1)$. The Eulerian number $\langle {n \atop k} \rangle$ counts permutations of $n$ elements with exactly $k$ ascents.

For our problem, we count permutations with exactly one ascent: $\langle {n \atop 1} \rangle$.

\subsection{Closed Form}

\begin{theorem}
$\displaystyle\left\langle {n \atop 1} \right\rangle = 2^n - n - 1$ for $n \ge 1$.
\end{theorem}

\begin{proof}
By induction using the recurrence $\langle {n \atop k} \rangle = (k+1)\langle {n-1 \atop k} \rangle + (n-k)\langle {n-1 \atop k-1} \rangle$.

Base: $\langle {1 \atop 1} \rangle = 0 = 2 - 1 - 1$.

Step: $\langle {n \atop 1} \rangle = 2\langle {n-1 \atop 1} \rangle + (n-1)\langle {n-1 \atop 0} \rangle = 2(2^{n-1} - n) + (n-1) = 2^n - n - 1$.
\end{proof}

\subsection{Main Formula}

Choosing $n$ distinct letters from 26 and counting arrangements with exactly one ascent:
\[
p(n) = \binom{26}{n} \cdot (2^n - n - 1)
\]

\subsection{Maximization}

Computing $p(n)$ for all $n$:

\begin{center}
\begin{tabular}{|c|r|}
\hline
$n$ & $p(n)$ \\
\hline
2 & 325 \\
5 & 1\,623\,414 \\
10 & 44\,186\,942\,025 \\
15 & 354\,025\,698\,756 \\
18 & 409\,511\,334\,375 \\
20 & 305\,498\,816\,850 \\
25 & 619\,225 \\
\hline
\end{tabular}
\end{center}

The maximum is achieved at $n = 18$.

\section{Result}

\[
\max_{1 \le n \le 26} p(n) = p(18) = \binom{26}{18}(2^{18} - 19) = \boxed{409511334375}
\]

\section{Answer}

\[
\boxed{409511334375}
\]

\end{document}
